% Folder 77, batch 1 (archivists' pages 1-20).
% Pass: Opus 5.5 (claude-opus-5-5), 2026-10-03 — first pass, unchecked against the pages by a human.
%
% Pages 1, 6 and 20 are covers inscribed in his hand; their words become
% section headings (no \page). Pages 2-5 are a black-ink run ("Digression
% quadratique": equivalent categories C_1(S)...C_10(S), reflections of a
% conic); pages 7-19 a blue-ink run on quadratic forms quasi-invariant under
% the dihedral group and conics circumscribed to regular polygons. Notation
% fixed for the batch: \zeta, \zeta^{-1} the eigenvalues of u, with
% \alpha = \zeta + \zeta^{-1}; \check{D}_i for the duals of the eigenlines
% D_i; \mathbb{D}_\infty for his outlined D; \tfrac12 kept as he writes it
% between two groups (presumably a semi-direct product); his underlined O
% (structure sheaf) is set as \mathcal{O}.
\documentclass[11pt,a4paper]{article}
\input{../preamble/grothendieck}

\folder{77}
\batch{1}
\pages{1}{20}
\foldertitle{Quadriques affines, extensions quadratiques : notes manuscrites (s.d.).}
\dating{[à partir de 1977]}
\watermark{Édition de démonstration}

\begin{document}

\section{Formes quadratiques de rg 2, droites projectives et revêtements quadr.}
\note{Titre repris de la couverture (p.~1 des archivistes), au crayon, de sa main ; le « et » devant « droites » est surchargé. Le feuillet ne porte rien d'autre.}


\page{2}
\emph{Digression quadratique}
\marginal{cf. aussi notes sur \uncertain{$\mathrm{GP}(1) \simeq \mathrm{SO}(3)$} et variantes}

$S$ schéma.

$C_1(S)$ catégorie des Modules quadratiques réguliers de rg~2 $(V, q)$.

$C_2(S)$ catégorie des \uncertain{quadruples} $(A, V, \xi)$, i.e.
\begin{itemize}
  \item $A$ Algèbre finie étale de rg~2 sur $S$,
  \item $V$ $A$-Module inversible
  \item et $\xi \in \Gamma N_{A/S}(V)^{*}$ \uncertain{section} de $N_{A/S}(V)$.
\end{itemize}

$C_3(S)$ catégorie des \struck{triples} \struck{\ill{}} $(U, T)$, $U$ une \uncertain{\struck{\ill{}}}-forme de $\mathbb{G}_m$, et $T$ \struck{\ill{}} un $U$-torseur.

$C_4(S)$ catégorie des couples $(X, \Delta)$, $X$ fibré en droites projectives sur $S$, $\Delta$ un diviseur relatif \add{étale} de rang~2.

$C_5(S)$ catégorie des triples $(P, X, D)$, $P$ fibré en \struck{droites} \add{plans} projectifs, $D$ \struck{droite} ss-fibré en droites projectives, $X$ conique lisse \uncertain{de $X$}\note{sic, semble-t-il : on attend « de $P$ ».}, $D$ et $X$ non tangents en chaque fibre.

$C_6(S)$ catégorie des couples $(E, C)$, $E$ fibré affine de rg~2 sur $S$, $C \subset E$ conique, projectivement lisse et \struck{partout} non parabolique en chaque fibre.

$C_7(S)$ \struck{triples} \add{quadruples} $(E, Q, \omega, V)$, $(E, Q, \omega)$ une forme \uncertain{spéciale}-orthogonale de rang~3, $V$ un sous-\struck{Module} fibré vectoriel de rg~2 \struck{tel} non isotrope [i.e. tel que $Q|V$ soit régulière].

$C_8(S)$ Catégorie des couples $(M, A)$, $M$ Alg. d'Azumaya de rg~4 sur $S$, $A$ sous-Algèbre commutative étale de rg~2.

\marginal{$C_9(S)$ : formes $G$ de \uncertain{$\mathrm{GP}(1)_S$}, avec \struck{\ill{}} un tore maximal \uncertain{de $G$} ; $C'_9(S)$ : formes $G$ de $\mathrm{SL}(2)_S$, avec un tore maximal \add{\uncertain{$W$}} de $G$.}
\note{Ces deux lignes sont écrites le long de la marge gauche, la feuille tournée.}

\page{3}
$C_{10}(S)$ Catégorie des torseurs sous $(\mathbb{Z}/2)_S \tfrac{1}{2}\, \mathbb{G}_{m,S}$.\note{Le signe « ½ » entre les deux groupes est le sien, ici et plus bas ; vraisemblablement un produit semi-direct.}

\struck{\ill{}} Toutes ces catégories sont canoniquement équivalentes…

\struck{Conditions équi} Pour \uncertain{$S'$} on a fixé (NB \uncertain{$S' = \mathrm{Spec}(A)$}…) les structures en question forment une \add{\ill{}} catégorie (claire sur $C_2$ et sur $C_3$). Si $T$ est un objet de celle-ci ($U$-torseur), on voit que $T^{\wedge 2}$ est loc. trivial au sens de Zariski.
\marginal{faux, \uncertain{cf.} sur $\mathbb{R}$, les formes quadratiques $\alpha(x^2+y^2)$ \ill{} $xy$ \ill{}}
On veut expliciter la question \struck{\ill{}} trivialisation de $T^{\wedge 2}$. Par la suite exacte de \struck{Kummer}
\[
  0 \to \mu_2 \to U \xrightarrow{\;2\;} U \to 0 ,
\]
on voit qu'il revient au même de se donner une restriction du groupe structural à $\mu_2$, ou une orbite \struck{\ill{}} de $T$ sous $\mu_2$, i.e. un morphisme de degré~2 $X \to X'$, où $X'$ est un fibré en droites \add{relatif} avec diviseur étale $\Delta'$, et une section \struck{\ill{}} de $X' \smallsetminus \Delta'$, $X \to X'$ induisant un morphisme \add{fini} étale $X \smallsetminus \Delta \to X' \smallsetminus \Delta'$. Ceci signifie aussi qu'on a une restriction des groupes structuraux
\[
  \mathbb{D}_S = \mathcal{O}(2)_S \quad \text{\uncertain{au} 4-groupe} \quad
  (\mathbb{Z}/2)_S \tfrac{1}{2}\, \mu_{2,S} \simeq (\mathbb{Z}/2)_S \times (\mu_2)_S
\]
qui opère sur $\mathcal{O}_S \times \mathcal{O}_S$ par symétrie et homothéties. Donc la catégorie de ces objets s'identifie à la cat. des triples $(A, L, \varepsilon)$, $A$ Algèbre finie étale de rg~2 sur $S$, $L$ module inv. sur $S$, $\varepsilon$ une trivialisation de $L^{\otimes 2}$ --- on prend $V = A \otimes_{O_S} L$, d'où $N_{A/O_S} V \simeq L^{\otimes 2}$, et on prend $\xi = \varepsilon$…

\marginal{NB L'action de $\mu_2$ sur $(X, \Delta)$ \ill{} qui \ill{} l'identité sur $\Delta$… (Cf \struck{\ill{}} \ill{} \ill{}…)}
\marginal{NB $X' = X/\mu_2$ \ill{} : \ill{} \ill{} \ill{} des $X' \smallsetminus \Delta'$ \ill{} $\mathbb{Z}/2\mathbb{Z}$ \ill{}}
\note{Trois notes marginales serrées, écrites en biais dans la marge gauche ; la première se rattache par un trait à « au sens de Zariski », les deux autres à la discussion de $X \to X'$. Seuls les fragments transcrits se lisent.}

\page{4}
\struck{$\tau_0\tau_2 = \tau_2\tau_0$ i.e.}

Données : \uncertain{un plan} de $X$ fibré en droites projectives \add{d'une Alg. de quat. $M$}) $\Rightarrow$ $C$ conique proj.

$\exists\, \tau_0, \tau_1, \tau_2 \in \mathrm{Aut}(X)$ réflexions (i.e. $\tau_i^2 = \mathrm{id}$, $\tau_i \neq 1$ en chaque fibre)
\[
  \Updownarrow
\]
$h_0, h_1, h_2 \in \Gamma(C \smallsetminus X)$
\[
  \Downarrow
\]
$L_0, L_1, L_2$ ortho (droites des pts fixes des $\tau_0^C, \tau_1^C, \tau_2^C$)

\emph{NB} Les $\tau_i$ se \uncertain{remontent} de \uncertain{façon} unique \struck{\ill{}} \add{à \ill{} près} en des \add{réflexions} \struck{\ill{}} de $M^{*}$.
\note{Passage très surchargé : plusieurs mots biffés et interlignés entre « unique » et « $M^{*}$ ».}

\emph{Hyp} :
\begin{itemize}
  \item[A)] $\tau_0\tau_2 = \tau_2\tau_0$, i.e. A') $h_0 \perp h_2$ ;
  \item[B)] $\tau_0\tau_1 \neq \tau_1\tau_0$, $\tau_1\tau_2 \neq \tau_2\tau_1$ sur chaque fibre, ou \uncertain{$\rho_{f}^2 \neq 1$, $\rho_{s}^2 \neq 1$}, i.e. B') $h_0 \not\perp h_1$, $h_1 \not\perp h_2$ --- (ou $h_0 \notin L_1$, $h_1 \notin L_2$ ---) ;
  \item[] \struck{C) \ill{}} \struck{C') $h_i \notin X$ ($i = 0,1,2$) sur chaque fibre}
  \item[C)] $\{h_0, h_1, h_2\}$ engendrent \uncertain{tout} $C$ en chaque fibre ($\Leftrightarrow$ $H_0, H_1, H_2$ indép. en chaque fibre) ;
  \item[D)] Considérons la droite \add{$H_1 \cap H_2$} de $\check{\mathbb{P}}(M) \supset C \simeq \check{\mathbb{P}}(\gamma M)$ formée des pts orth. à $h_1$, \struck{\ill{}}, $h_2$, i.e. fixes par $\tau_1, \tau_2$. \add{Les conditions précédentes \ill{} qu'elle \ill{}} \struck{\ill{} qu'elle \ill{}} elle n'est pas contenue dans la \uncertain{quadrique} \ill{} de $\check{\mathbb{P}}(M)$ définie par ces \uncertain{deux} \ill{}). On exige qu'elle ne soit pas \add{en chaque fibre} tangente, ce qui signifie aussi que \uncertain{$\beta \neq 2$} inv., on \ill{} a) Aux pts \add{de $S$} de car. \add{impaire} $p > 0$, on veut que l'ordre de $\tau_1\tau_2$ soit $\neq p$ (car \ill{} \ill{} \struck{\ill{}} \add{deux pts} \struck{\ill{}}) \ill{} aux pts de car. $> 0$ \struck{\ill{}}, \ill{} si $\tau_1\tau_2$ d'ordre fini \ill{} \add{\ill{}} aux pts de car. nulle.
\end{itemize}
\note{La fin de D), en bas de page, est très biffée et surchargée ; seuls les mots transcrits se lisent avec quelque assurance.}

\marginal{ce qui, pour $\tau_i^2 = \mathrm{id}$ ($i = 0,1,2$), implique que \uncertain{les} $\tau_i$ sont des \uncertain{réflexions} (car $\tau_i^2 \neq 1$ sur chaque fibre)}
\marginal{ceci \uncertain{implique} aussi \ill{} que \ill{} en chaque fibre}
\marginal{\uncertain{dire} que $h_i \notin X$ ($i = 0,1,2$) --- comment \uncertain{l'exprimer} en termes des opérations $\tau_i$ sur $X$ ?}
\marginal{il suffit \ill{} \uncertain{seulement} \ill{} des $S$ \ill{} car. $p > 0$ !}
\note{Notes de la marge gauche, reliées par des flèches et un trait courbe aux conditions B) et C). Dans la première, il écrit $\tau_i^2 \neq 1$ là où le corps de page a $\tau_i \neq 1$.}

\page{5}
\[
  X \longleftrightarrow M \longleftrightarrow (X \subset C), \qquad C = \check{\mathbb{P}}(\gamma M)
\]

\begin{itemize}
  \item réflexions de \struck{$X$} \add{$X$} ;
  \item[$\updownarrow$] réflexions de $C$ stables par $X$ $\longleftrightarrow$ réflexions de $V$ orth. \quad (\uncertain{$V \simeq \check{\gamma}M$})
  \item[$\updownarrow$] pts de $C \smallsetminus X$
  \item[$\updownarrow$] droites de \struck{$M$} \add{$\gamma M$} non \uncertain{isotropes} \qquad $M^{*}$ : \uncertain{Tr nulle}
  \item[$\updownarrow$] div\add{iseurs étales} de deg~2 sur $X$
  \item[$\updownarrow$] \struck{droites} droites de $C$ non tgtes à $X$
  \item[$\updownarrow$] plans de $V$ (i.e. de $\gamma M$) non tangents \uncertain{au cône} $q = 0$
  \item[$\updownarrow$] sous-algèbres étales de rg~2 de $M$
  \item[$\updownarrow$] Éléments $\sigma$ de $M$ (tels que \struck{\ill{}}, $\sigma^2 = 1$, $\sigma \neq \pm 1$ en chaque fibre \add{\uncertain{quotients des réflexions}, i.e.} \struck{\ill{}}  --- conditions \struck{\ill{}} \uncertain{rien. pour la} \uncertain{remplacer par} ($\det \sigma = -1$))
\end{itemize}
\note{Il dispose ces objets en colonne, reliés par des flèches doubles verticales ; une accolade en marge gauche porte « car.~$\neq 2$ » en deux endroits, et un renvoi : « car.~$\neq 2$, \uncertain{sinon} il faut \ill{} \ill{} ».}

L'action (via transport de structures) de la réflexion $\tau$ sur $M$, en termes de $\sigma$, est $u \mapsto \sigma u \sigma^{-1}$, mais ce n'est pas une réflexion de $M$ ! (val. propres $1, -1$ avec mult.~2 l'une et l'autre). Pour avoir une réflexion il faut \add{composer} \struck{\ill{}}… $u \mapsto u^{*} = \mathrm{Tr}\, u . 1 - u$ !!!

On \uncertain{a} :
\[
  \tilde{\sigma}(u) = \sigma u^{*} \sigma = (\sigma u \sigma)^{*} = u - \mathrm{Tr}(u \sigma)\, \sigma
\]
\note{La dernière formule est encadrée.}

\section{Formes quadratiques inv. et quasi-inv. sous $D_\infty$, $D_n$}
\note{Titre repris de la couverture (p.~6 des archivistes), au crayon, de sa main ; elle porte aussi la mention « (13~p) ». Les pp.~7--19 de ce lot font treize feuillets.}

\page{7}
\emph{Formes quadratiques quasi-invariantes sous le groupe diédral.}

\emph{Cas I} \quad $\alpha \neq \pm 2$ fibre par fibre \add{($\alpha^2 - 4$ inv.)}, de sorte que la repr. \struck{affine} \add{projective} (\struck{\ill{}} affine) naturelle de $\mathbb{D}_\infty = \mathrm{Aut}(\uncertain{\Omega})$ se restreint en une représentation vectorielle où
\[
  (1) \qquad
  \begin{aligned}
    \sigma(x,y) &= (y, x), \qquad \sigma'(x,y) = (\zeta y, \zeta^{-1} x), \\
    u(x,y) &= (\zeta x, \zeta^{-1} y) .
  \end{aligned}
\]
La décomposition $E = D_1 \oplus D_2 \oplus D_0$ (correspondant aux val. propres $\zeta, \zeta^{-1}$ de $u$, $\zeta \neq \zeta^{-1}$ fibre par fibre) définit une décomposition analogue de
\[
  (2) \qquad
  \begin{aligned}
    \mathrm{Sym}^2(\check{E}) \simeq{}& \check{D}_1^{\otimes 2} \oplus \check{D}_2^{\otimes 2} \oplus \check{D}_0^{\otimes 2} \\
    &\oplus \check{D}_1 \otimes \check{D}_2 \oplus \check{D}_1 \otimes \check{D}_0 + \check{D}_2 \otimes \check{D}_0
  \end{aligned}
\]
correspondant aux valeurs propres
\[
  (3) \qquad \zeta^{2},\ \zeta^{-2},\ 1,\ 1,\ \zeta^{-1},\ \uncertain{\zeta}
\]
de $u$.\note{Il écrit $\zeta^{2}, \zeta^{-2}$ pour $\check{D}_1^{\otimes 2}, \check{D}_2^{\otimes 2}$ ; sur le dual on attendrait $\zeta^{-2}, \zeta^{2}$. Transcrit tel quel.}

Supposons d'abord les \add{5} val. propres $1, \zeta, \zeta^{-1}, \zeta^{2}, \zeta^{-2}$ distinctes deux à deux fibre par fibre (ce qui signifie $\zeta^{3} \neq 1$, $\zeta^{4} \neq 1$ fibre par fibre, en plus de $\zeta \neq \pm 1$ \add{fibre par fibre, i.e. $(\zeta^2 - 1)$ inv.,} exprimant $\alpha \neq \pm 2$ fibre par fibre…). Alors les droites propres pour $u$ sont \struck{\ill{}} d'une part $\check{D}_1^{\otimes 2}$, $\check{D}_2^{\otimes 2}$, $\check{D}_1 \otimes \check{D}_0$, $\check{D}_2 \otimes \check{D}_0$, d'autre part les droites contenues \add{\uncertain{(ens)}} dans $\check{D}_0^{\otimes 2} \oplus (\check{D}_1 \otimes \check{D}_2)$. Or les 4 droites précédentes sont non inv. par $\sigma$ (qui échange $\check{D}_1$ et $\check{D}_2$, donc $\check{D}_1^{\otimes 2}$ et $\check{D}_2^{\otimes 2}$, $\check{D}_1 \otimes \check{D}_0$ et $\check{D}_2 \otimes \check{D}_0$), d'autre part $\sigma$ opère trivialement dans $\check{D}_0^{\otimes 2} \oplus (\check{D}_1 \otimes \check{D}_2)$. Donc on trouve, \add{\uncertain{en somme}} \struck{les \ill{}}

\page{8}
comme seules droites invariantes \add{par $\mathbb{D}_\infty$} dans $\mathrm{Sym}^2(\check{E})$, \struck{les} celles qui sont formées des formes quadratiques \emph{invariantes}.

Si on permet que $\zeta^{3} = 1$ sur certaines fibres, \struck{\ill{}} mais pas $\zeta^{4} = 1$, alors on trouve que les droites inv. dans $\mathrm{Sym}^2(\check{E})$ par $u$ sont contenues \add{les} dans $V = \check{D}_1^{\otimes 2} \oplus \check{D}_2$ ou $W = \check{D}_2^{\otimes 2} \oplus \check{D}_1$, ou $\check{D}_0^{\otimes 2} \oplus \check{D}_1 \otimes \check{D}_2$, or $\sigma$ échange $V$ et $W$, donc il n'y a pas de droites inv. par $\mathbb{D}$ ni dans $V$, ni dans $W$, et on \struck{\ill{}} \add{obtient} la \uncertain{même} conclusion.
\marginal{i.e. $\alpha = -1$ sur certaines fibres}

Mais si on permet $\zeta^{4} = 1$ sur certaines fibres, que se passe-t-il ? Supposons donc carrément
\[
  (4) \qquad \zeta^{4} = 1 \quad \text{i.e.} \quad \zeta^{2} = -1, \quad \text{ou encore } \alpha = 0
\]
(cas du carré) --- car (4) s'écrit
\[
  (\zeta^{2} - 1)(\zeta^{2} + 1) = 0
\]
et comme par hyp. $\zeta^{2} - 1$ est inv., cela donne $\zeta^{2} + 1 = 0$ i.e. $\zeta + \zeta^{-1} = 0$ i.e. $\alpha = 0$…

Alors les valeurs propres distinctes parmi (3) \struck{\ill{}} sont $\zeta^{2} = \zeta^{-2} = -1$, $1$, $\zeta, \zeta^{-1}$, correspondant aux sous-espaces propres $\check{D}_1^{\otimes 2} \oplus \check{D}_2^{\otimes 2}$, $\check{D}_0^{\otimes 2} \oplus (\check{D}_1 \otimes \check{D}_2)$, $\check{D}_1$, $\check{D}_2$. Les droites propres

\page{9}
sous $u$ sont celles contenues dans un de ces sous-espaces. Si on veut qu'elles soient aussi inv. sous $\sigma$, il ne reste plus que les droites des deux premiers sous-espaces. Or on est en car. $\neq 2$ (car $\zeta^{2} = -1$, mais $\zeta \neq \pm 1$ fibre par fibre) et $\sigma$ opère sur $\check{D}_1^{\otimes 2} \oplus \check{D}_2^{\otimes 2}$ en permutant $\check{D}_1^{\otimes 2}$ et $\check{D}_2^{\otimes 2}$… Cela fait donc 2 droites propres canoniques de l'espace qu'il engendre, correspondant aux val. propres $+1$ et $-1$ resp. pour $\sigma$, les val. propres pour $u$ étant $-1, -1$, celles pour $\sigma' \add{= u\sigma}$ \struck{\ill{}} étant $-1, 1$.

Si on apprécie le carré en coordonnées $(x,y)$ avec sommets $(\pm 1, \pm 1)$ \struck{\ill{}} rangés dans l'ordre
\[
  (1,1),\ (-1,1),\ (-1,-1),\ (1,-1)
\]
\drawing{Petit carré quadrillé par les deux axes, sommets marqués $s_0$ (en haut à droite), $s_1$ (en haut à gauche), $s_2$ (en bas à gauche), $s_3$ (en bas à droite).}
\marginal{NB car $\neq 2$}

Les formes quadratiques en $(x,y,z)$ correspondantes sont $x^2 - y^2$ \add{$= (x-y)(x+y)$} et $xy$ (NB $\sigma(x,y) = (y,x)$, $\sigma'(x,y) = (-x, y)$)
\struck{elles sont toutes deux \add{\ill{}} dégénérées} $\{$ La première prend 2 valeurs constantes $\neq 0$ sur les sommets du carré, la seconde prend la valeur $+1$ sur \add{un couple de} \struck{deux} sommets opposés, la valeur $-1$ sur le couple des deux autres sommets… $\}$ Seule la première définit une conique \uncertain{non} \ill{} \ill{} par les sommets du carré \uncertain{passant} \ill{}
\note{La dernière ligne, au ras du bord inférieur, est en partie coupée ; elle se poursuit en haut de la p.~10.}

\page{10}
\uncertain{conique $x^2 - y^2 = 0$}, mais celle-ci a l'origine comme pt singulier \add{\uncertain{i.e.} réunion des deux droites diagonales…}\note{Cette ligne, et l'interligne au-dessus de « On voit », sont ajoutées en tête de page, très serrées.}
\marginal{bref !}
On voit \emph{donc} que si $\zeta^{4} - 1$ est inv. i.e. $4(\alpha^2 - 4)$ inv., donc \struck{\ill{}} les formes quadratiques sur $E$ quasi-invariantes par $\mathbb{D}$ et \emph{non dég.} sur chaque fibre sont en fait invariantes par $\mathbb{D}$, et c'est une combinaison linéaire des formes $z^2$ et $xy$ (avec les coordonnées choisies), \struck{\ill{}} où on a dû d'abord faire un changement de base quadratique avec
\[
  S' = \mathrm{Spec}\, \mathcal{O}_S[T]/(T^2 - \alpha T + 1)\ \ldots)
\]
\[
  (5) \qquad Q(x,y,z) = a z^2 + b xy
\]
dont la \struck{déterminant} matrice
\[
  \begin{pmatrix} 0 & b & 0 \\ b & 0 & 0 \\ 0 & 0 & 2a \end{pmatrix}
\]
a le déterminant
\[
  (6) \qquad \delta(Q) = 2ab^2
\]
donc le déterminant divisé
\[
  (7) \qquad \delta'(Q) = ab^2
\]
\uncertain{Quasi}-\struck{\ill{}} \add{inv.} ssi $a, b$ le sont. La condition que $Q$ s'annule sur les sommets de $P$ s'écrit alors
\[
  (8) \qquad a + b = 0
\]
donc $\exists !$ \struck{\ill{}} conique circonscrite à $P$ invariante sous $\mathbb{D}_\infty$, \struck{\ill{}} d'équation
\[
  (9) \qquad z^2 - xy = 0
\]
On en conclut, en \uncertain{condensé}, \uncertain{affirm.}, $\exists !$

\page{11}
\struck{fonction quadratique} \add{polyn. de degré $\leq 2$} \struck{sur} \add{\uncertain{inscrite} \uncertain{puis}} $V$ invariante par $\mathbb{D}_\infty$ et \struck{\ill{}} nulle sur les sommets du polygone $P$, savoir
\[
  xy - 1
\]
[en somme, $\exists !$ forme quadratique sur $V$ invariante par $\mathbb{D}_\infty$ égale à $1$ sur les sommets du polygone $P$, savoir $xy$…].

\emph{NB} Il y a « en général » une et une seule conique qui passe par cinq points, (NB cinq équations lin. homogènes en 6 coeff. homogènes…), donc on s'attend à ce que, lorsque $P$ a au moins cinq pts fibre par fibre (i.e. $\alpha \neq 0, -1$ fibre par fibre pour exclure triangles et carrés, donc au total
\[
  \alpha(\alpha^2 - 4)(\alpha + 1) \quad \text{inversible}
\]
qu'il existe \add{une unique} \struck{\ill{}} conique qui passe par les sommets de $P$. Écrivons la conique sous forme non homogène comme
\[
  Q(x,y) = a x^2 + 2b\, xy + c y^2 + u x + v y + \lambda
\]
et prenons pour $(x,y)$ les cinq points $(\zeta^{i}, \zeta^{-i})$, $-2 \leq i \leq +2$, on trouve

\page{12}
les conditions sur $a, b, c, u, v, \lambda$ s'expriment par un syst. de 5 équations linéaires, de matrice
\[
  \begin{pmatrix}
    1 & 1 & 1 & 1 & 1 & 1 \\
    \zeta^{2} & 1 & \zeta^{-2} & \zeta & \zeta^{-1} & 1 \\
    \zeta^{-2} & 1 & \zeta^{2} & \zeta^{-1} & \zeta & 1 \\
    \zeta^{4} & 1 & \zeta^{-4} & \zeta^{2} & \zeta^{-2} & 1 \\
    \zeta^{-4} & 1 & \zeta^{4} & \zeta^{-2} & \zeta^{2} & 1
  \end{pmatrix}
\]
et il faut vérifier si \struck{\ill{}} \add{cette matrice ($5 \times 6$)} est de rang~5 fibre par fibre, ce qui équivaut (vu qu'il y a deux colonnes identiques formées de 1) que le déterminant \struck{\ill{}}
\[
  \det
  \begin{pmatrix}
    1 & 1 & 1 & 1 & 1 \\
    \zeta^{2} & \zeta^{-2} & \zeta & \zeta^{-1} & 1 \\
    \zeta^{-2} & \zeta^{2} & \zeta^{-1} & \zeta & 1 \\
    \zeta^{4} & \zeta^{-4} & \zeta^{2} & \zeta^{-2} & 1 \\
    \zeta^{-4} & \zeta^{4} & \zeta^{-2} & \zeta^{2} & 1
  \end{pmatrix}
  = \Delta(\zeta), \qquad \text{où } \Delta \in \mathbb{Z}[T, T^{-1}]
\]
est inversible. De la façon dont on obtient ce déterminant, il \struck{semble} est clair qu'il est nul si \struck{\ill{}} $\zeta^{3} = 1$, ou $\zeta^{4} = 1$, \add{\uncertain{faut} une des six \uncertain{premières} \uncertain{lettres} \uncertain{simplifiées} \uncertain{dans} $Q$) \uncertain{du} \uncertain{tableau}}, i.e. si
\[
  \frac{(\zeta^{4} - 1)(\zeta^{3} - 1)}{\zeta - 1} \in \mathbb{Z}[\zeta]
\]
\struck{et \ill{} par $(\zeta^2+1)$}
\marginal{\uncertain{manœuvre} Vandermonde des colonnes, \uncertain{bien} \uncertain{voir} matrices symétriques}
on trouve que $\Delta(T)$ est divisible par
\[
  \frac{(\zeta^{2}-1)(\zeta^{3}-1)}{\zeta - 1} = (\zeta^{2}+1)(\zeta^{2}-1)(1 + \zeta + \zeta^{2}), \quad \text{et} \equiv \text{par}
\]
\[
  (\zeta^{2}+1)(\zeta^{2}-1)^{2}(1 + \zeta + \zeta^{2}) = \alpha(\alpha^{2} - 4)(\alpha + 1)
\]
\uncertain{Étudier} la multiplicité dudit déterminant ?
\note{Dans la première fraction de la ligne avant-dernière il écrit $(\zeta^{2}-1)$ au numérateur, contre $(\zeta^{4}-1)$ trois lignes plus haut. La dernière identité, prise à la lettre, demande de diviser par une puissance de $\zeta$. Transcrit tel quel.}

\page{13}
\ill{} \add{\uncertain{fonction} pol. de degré $\leq 2$} $q$ (qui s'annule sur cinq \uncertain{sommets}\note{Le mot « sommets », en fin de ligne, est d'une autre encre (brune) et en écriture détachée ; le début de la ligne est couvert par une tache.}) consécutifs s'annule sur tous les sommets (si $\alpha(\alpha+1)$ inversible) et est en fait un multiple de la forme \struck{quadratique} invariante principale. Sur l'espace des translations, la forme quad. principale est
\[
  (3) \qquad \text{\struck{$q(x,y)$}}\ \ q_0(x,y) = x^2 + y^2 - \alpha xy
\]
\struck{qui est dég.} \add{dont le} discriminant est
\[
  (4) \qquad \delta_0 = \alpha^2 - 4 = (\alpha - 2)(\alpha + 2)
\]
inv. ssi $\alpha \neq \pm 2$ fibre par fibre i.e. ssi on est dans le cas~I.

Quant à $q_0$ elle-même, écrite comme forme en 3 variables $x, y, z$
\[
  (2\,\text{bis}) \qquad q_0(x,y,z) = x^2 + y^2 \ \text{\struck{$+ \ill{}$}} + 0 z^2 - yz - \text{\struck{\ill{}}}\, xz - \alpha xy
\]
de matrice \add{\uncertain{adjointe}}
\[
  (5) \qquad
  \begin{pmatrix} 2 & -\alpha & -1 \\ -\alpha & 2 & -1 \\ -1 & -1 & 0 \end{pmatrix}
\]
de discriminant
\[
  (6) \qquad \delta = \text{\struck{\ill{}}} -2\alpha - 2 - 2 = -2(\alpha + 2)
\]
de discriminant divisé
\[
  (6') \qquad \delta' = -(\alpha + 2)
\]
qui est nul ssi $\alpha = -2$, inv. ssi $\alpha \neq -2$ sur chaque fibre (i.e. si on n'est pas dans le cas~III sur aucune fibre…)
\note{Les numéros d'équations (3), (4), (2 bis), (5), (6) ne prolongent pas ceux des pp.~7--12, qui vont jusqu'à (9) : ce feuillet semble appartenir à une autre rédaction du même calcul.}

\page{14}
Les deux pts isotropes (confondus \struck{\ill{}} en chaque fibre au-dessus de $s$ ssi $\alpha_s = \pm 2$…) correspondent aux \struck{solutions} directions de pente $\zeta = y/x$ satisfaisant
\[
  (7) \qquad \zeta^2 - \alpha\zeta + 1 = 0 \quad \text{i.e.} \quad \zeta + \zeta^{-1} = \alpha
\]
Étudions la situation où on se donne (en plus du polygone régulier) une telle direction isotrope --- ici le \struck{schéma} « modulaire » est donc \add{la racine inversible} $\zeta$, plus fin que $\alpha$ [et pour les « \struck{\ill{}} $n$-polygones » on aura $\Phi_n(\zeta) = 0$ \add{$\to$ pol. cyclotomique} + conditions de caractéristique…]. Donc la forme quadratique $q_0$ s'écrit
\[
  (8) \qquad q_0(x,y) = (x - \zeta y)(x - \zeta^{-1} y)
\]
et $q(x,y,z)$ s'écrit
\[
  (9) \qquad q(x,y,z) = (x - \zeta y)(x - \zeta^{-1} y) - (x + y) z
\]
\note{La numérotation repart ici de (7), à la suite de celle de la p.~13.}
Utilisons \struck{\ill{}} \add{\struck{correspondant au point}} la \struck{\ill{}} \add{droite $D_\zeta$ de pente $\zeta$ et} (le pt isotrope $d_\zeta \in C$) pour identifier la conique $C$ : $q(x,y,z) = 0$ à $\mathbb{P}(E/D_\zeta)$ --- on trouve…

\drawing{En marge gauche : une ellipse, un point marqué $d_\zeta$ sur elle, et deux droites issues de ce point, l'une tangente, l'autre sécante.}

Un repère projectif de $\mathbb{P}(E/D_\zeta)$ \uncertain{grâce} aux trois plans suivants contenant $D_\zeta$ i.e. aux droites projectives de $\check{\mathbb{P}}(E)$ contenant $d_\zeta$ :
\begin{itemize}
  \item[$\circ$)] \struck{le plan qui} \uncertain{la plane} \struck{\ill{}} \ill{} : la droite tangente : $d_\zeta$ : $C$
\end{itemize}
\note{La fin de la page est surchargée ; l'énumération des trois plans se poursuit à la page suivante.}

\page{15}
\[
  \left\{
  \begin{aligned}
    f_1 &= (1, \zeta, 0) = e_1 + \zeta e_2 \\
    f_2 &= (1, 0, 1 - \zeta) = e_1 + (1 - \zeta) e_3 \\
    f_3 &= (0, 0, 1) = e_3
  \end{aligned}
  \right.
\]
\struck{$f_3 = (0, \zeta, \zeta - \ill{}) = e_2 + (\zeta - 1) e_3$ --- vecteur directeur de la 2e droite isotrope}
\drawing{Petit croquis en perspective : trois vecteurs $f_1$, $f_2$, $f_3$ issus d'une origine, un segment marqué $1 - \zeta$, et un plan marqué $H_\zeta$.}

\emph{NB}
\[
  \begin{aligned}
    e_1 &= f_2 - (1 - \zeta) f_3 = f_2 + (\zeta - 1) f_3 \\
    e_2 &= \zeta^{-1}(f_1 - e_1) = \zeta^{-1} f_1 - \zeta^{-1} f_2 + (1 - \zeta^{-1}) f_3 \\
    e_3 &= f_3
  \end{aligned}
\]
\struck{$\ill{} = \zeta^{-1} f_1 - \zeta^{-1} f_2 + (\zeta - \zeta^{-1}) f_3$}
\note{Dans ces lignes il a corrigé plusieurs indices par surcharge ($f_2$ / $f_3$) ; on transcrit la leçon finale. La ligne biffée en tête de colonne droite est barrée de traits obliques.}

Le plan \add{$(e_1, e_3)$ est invariant} engendré par $(f_1, e_1)$ (qui correspond à la droite isotrope $D_{\zeta^{-1}} = D'_\zeta$…), qui correspond dans le plan $(f_2, f_3)$ --- \uncertain{si} $\zeta - 1$ inv. --- à la droite de pente (par rapport à $f_3, f_2$) $\dfrac{1}{\zeta - 1}$

\struck{Donnée des Polygones épinglés \add{par $D_\zeta$} du type $\neq$ III \add{\uncertain{i.e. réguliers}} \uncertain{équivalent} + droite isotrope $\Longleftrightarrow$ donnée de $\zeta \in \Gamma(S, \mathcal{O}_S)^{*}$ $\Longleftrightarrow$ ce qui suit :}
\begin{itemize}
  \item[\struck{1°)}] \struck{fibré en \struck{droites affines \ill{}} i.e. torseurs sous \add{$F$} fibré vectoriel inv. \add{de rang 1} \struck{\ill{}} i.e. extension \uncertain{via} $e_\zeta$ par un $L$ \quad $0 \to L \to F \to \mathcal{O}_S \to 0$}
  \item[\struck{2°)}] \struck{Une section de $\mathbb{P}(F)$ i.e. un Module inv. $L'$ et un hom. $F \to L'$}
\end{itemize}
\note{Tout ce bloc, depuis « Donnée des Polygones », est barré de deux longs traits obliques, qui s'arrêtent avant « Avec ceci ».}

Avec ceci, on construit le schéma en groupes

\page{16}
\struck{des automorphismes de $P$ qui laissent invariants les diviseurs \struck{sections} relatifs de degré~2 qui sont définis par les 2 sections de $P$, \emph{et} induisent l'identité dessus --- qui est un schéma en groupes $G$ lisse sur $S$, de dim relative~1, dont la comp. \struck{neutre} $G^0$ est commutative}

\struck{3°)}
\note{Ces lignes, suite de la phrase ouverte au bas de la p.~15, sont barrées de deux traits obliques ; le reste du feuillet est blanc.}

\page{17}
Quand aux formes quadratiques invariantes \add{\uncertain{quasi-inv.}} sur l'espace des translations $\check{V}$ (toujours sous l'hyp. $\alpha^2 - 4$ inv.) elles sont comprises \uncertain{parmi} les quasi-inv. sur $\check{E}$ (dont \struck{\ill{}} \struck{\ill{}} \add{\struck{\ill{}} est un quotient}, puisque $\alpha \neq 2$ fibre par fibre…), on trouve donc, si $\alpha$ inv., les multiples seulement de la forme $xy$ (qui est non dégénérée et donne lieu à la conique circonscrite à $P$), et si $\alpha = 0$ i.e. \struck{\ill{}} $\zeta^2 = -1$ \add{(donc 2 inv. sur $S$)} on trouve comme dessus \struck{les} multiples de $x^2 + y^2$ et de $x^2 - y^2$, dont la \add{première prend les valeurs} \struck{\ill{} \ill{} \ill{}} \ill{} \uncertain{tels} 2 et $\zeta^2 + \zeta^{-2} = \alpha^2 - 2 = -2$ sur les sommets de $P$ \add{(donc deux valeurs distinctes)}, et la deuxième y prend 2 valeurs $0$, mais (comme nous avons remarqué) définit une conique circonscrite dégénérée….

Il reste à étudier la question des formes quadratiques \add{quasi} \struck{invariantes}, d'abord pour $\alpha = 2$ et $\alpha = -2$, puis éventuellement pour $\alpha$ « variable » général….

\page{18}
\emph{Coniques circonscrites à un polygone régulier.}
\note{Le titre et la ligne suivante sont repassés à l'encre brune sur certains mots (« polygone régulier », « universelles »), comme le mot « sommets » de la p.~13.}

Faisons les calculs en coordonnées universelles
\[
  (1) \qquad
  \left\{
  \begin{aligned}
    s_0 &= (0,0) \\
    s_1 &= (1,0) \\
    s_{-1} &= (0,1) \\
    s_2 &= (1 + \alpha, 1) \\
    s_{-2} &= (1, 1 + \alpha)
  \end{aligned}
  \right.
\]
Écrivons une fonction quadratique générale en coord. affines
\[
  q(x,y) = a x^2 + b xy + c y^2 + u x + v y + d
\]
\struck{\ill{}} Écrivons que la \add{\uncertain{$S$ \ill{}}} conique qu'elle définit \add{\uncertain{$-2 \leq i \leq +2$}} passe par les $s_i$, on trouve
\[
  \left.
  \begin{aligned}
    q(s_0) &= d = 0 \\
    q(s_1) &= a + u = 0 && \text{i.e. } u = -a \\
    q(s_{-1}) &= c + v = 0 && \phantom{\text{i.e. }} v = -b
  \end{aligned}
  \right\}
  \quad q(x,y) = a(x^2 - x) + c(y^2 - y) + bxy
\]
\note{Dans la troisième ligne, le coefficient de $q(s_{-1})$ est surchargé (un $c$ et un $b$ l'un sur l'autre) et la ligne conclut « $v = -b$ » ; la forme résumée à droite porte $c(y^2-y)$, ce qui demande $v = -c$. Transcrit tel quel.}
\[
  \left\{
  \begin{aligned}
    q(s_2) &= a(1+\alpha)\alpha + b(1+\alpha) = (1+\alpha)(\alpha a + b) = 0 \\
    q(s_{-2}) &= c(1+\alpha)\alpha + b(1+\alpha) = (1+\alpha)(\alpha c + b) = 0
  \end{aligned}
  \right.
\]
Si $\alpha - 1$ inv. (i.e. on n'a de triangle sur aucune fibre) la condition s'écrit
\[
  b = -\alpha a \ \text{\struck{\ill{}}} = -\alpha c
\]
ou encore et si $\alpha$ inversible (i.e. on n'a pas de carrés sur aucune fibre)
\[
  a = c, \quad b = -\alpha a,
\]
\add{i.e.} $q$ doit être multiple \add{$a q_0$} de la forme
\[
  (2) \qquad q_0(x,y) = (x^2 - x) + (y^2 - y) - \alpha xy
\]
Comme celle-ci est invariante sous $G$, \uncertain{on}
\marginal{NB C'est la conique universelle qu'il \ill{} si \ill{} --- vérif $q(s_2, s_{-2}) = \ill{}$ (\uncertain{erreur} \ill{})…}
\note{Note écrite en biais dans le coin inférieur gauche, reliée par des traits aux lignes « aucune fibre » et (2) ; très peu lisible. Dans « Si $\alpha - 1$ inv. », le facteur en jeu est $(1+\alpha)$ ; transcrit tel quel.}

\page{19}
\note{Ce feuillet prolonge l'énumération des plans contenant $D_\zeta$ commencée au bas de la p.~14 et les vecteurs $f_i$, $e_i$ de la p.~15, plutôt que la p.~18.}
\struck{b) celui qui correspond au plan $(x,y)$}

Plan tangent : $C$ le long de $D_\zeta$. Son équation est (utilisons la matrice adjointe \struck{\ill{}})
\[
  Ax + By + Cz = 0 \qquad \text{\struck{$2x - \alpha y =$}}
\]
où $A, B, C$ sont donnés par $\Delta\cdot(1, \zeta, 0)$ i.e.
\[
  \begin{aligned}
    A &= 2 - \alpha\zeta = 1 - \zeta^2 \ \text{\struck{$= \ill{}$}} = (1 - \zeta)(1 + \zeta), \\
    B &= -\alpha + 2\zeta = \zeta - \zeta^{-1} = \zeta^{-1}(1 - \zeta^2) = -\zeta^{-1}(1 - \zeta)(1 + \zeta), \\
    C &= -(1 + \zeta)
  \end{aligned}
\]
\note{Dans $B$, les deux dernières expressions sont de signes opposés ; transcrit tel quel.}

On peut (comme $1 + \zeta$ \struck{\ill{}} inv. i.e. $\zeta \neq -1$ fibre par fibre i.e. $\alpha \neq -2$ fibre par fibre, par hyp.) diviser par $(1 + \zeta)$, on trouve l'équation
\[
  H_\zeta : \quad
  \left\{
  \begin{aligned}
    &(1 - \zeta)x - \zeta^{-1}(1 - \zeta) y - z = 0 \quad \text{i.e.} \\
    &z = (1 - \zeta)(x - \zeta^{-1} y)
  \end{aligned}
  \right.
\]
\[
  D_\zeta \subset H_\zeta \subset E, \qquad H_\zeta \simeq V
\]
\note{Sous $H_\zeta$, un « $\simeq$ » vertical renvoie à $V$.}
Ainsi $C$ s'identifie à la \add{droite projective de la} droite affine associée à l'extension de $E/H_\zeta$ \add{$(\mathcal{O}_S e_3)$} par $H_\zeta/D_\zeta \simeq V/D_\zeta \simeq$. Comme l'ext. des 2 \add{est une} \uncertain{scindée} \struck{\ill{}} \add{\uncertain{extension} \uncertain{affine} \uncertain{est}} \uncertain{vectorielle}, isomorphe au fibré vectoriel cov. associé :
\[
  (E/H_\zeta)^{\otimes -1} \otimes H_\zeta/D_\zeta, \qquad
  E/H_\zeta \simeq \mathcal{O}_S\, e_3, \quad H_\zeta/D_\zeta \simeq \mathcal{O}_S\, e_1
\]
\drawing{En bas à gauche : deux axes $e_1$ (horizontal) et $e_3$ (vertical), un vecteur $e_2$ et la droite $D_\zeta = \mathcal{O}_S(1,\zeta)$ issus de l'origine.}
\note{La phrase sur l'extension, très surchargée d'interlignes, n'est lue qu'en partie ; l'argument se poursuit au-delà de ce lot.}

\section{Extensions quadratiques}
\note{Titre repris de la couverture (p.~20 des archivistes), à l'encre, de sa main ; elle porte aussi, au crayon, la mention « (43~p) ». La liasse qu'elle annonce commence au lot suivant.}

\end{document}
